
\documentclass[12pt]{article}
\usepackage[margin=1in]{geometry}
\usepackage{amsmath,amsthm,amssymb,scrextend}
\setlength{\parindent}{0pt}
\begin{document}
\newcommand{\ds}{\displaystyle}

\textbf{(a)(i)}
$$f(12)=12\left(1-\frac12\right)\left(1-\frac13\right)\\=\boxed{4}$$

$$f(180)=180\left(1-\frac12\right)\left(1-\frac13\right)\left(1-\frac15\right)=\boxed{48}$$

\bigskip\hrule\bigskip

\textbf{(a)(ii)}

Let $\quad\displaystyle  N=p_1^{\alpha_1}p_2^{\alpha_2}\cdots p_k^{\alpha_k}$, where $p_1,\ldots,p_k$ are distinct primes and $\alpha_1,\ldots,\alpha_k$ are positive integers.
\begin{align*}
f(N)
&=p_1^{\alpha_1}\cdots p_k^{\alpha_k}
\left(1-\frac1{p_1}\right)\cdots
\left(1-\frac1{p_k}\right)\\
&=p_1^{\alpha_1-1}(p_1-1)\cdots
p_k^{\alpha_k-1}(p_k-1)
\end{align*}

Thus $\displaystyle f(N)\in\mathbb Z$

\bigskip\hrule\bigskip

\textbf{(b)(i)} $\quad \boxed{\text{False}}$
$$f(3)=2,\quad f(9)=6, \quad f(27)=18$$
$$f(3)f(9)=12\neq18=f(27)$$


\bigskip\hrule\bigskip

\textbf{(b)(ii)} $\quad \boxed{\text{True}}$
\begin{align*}
f(p)f(q)
&=p\left(1-\frac1p\right)
q\left(1-\frac1q\right)\\
&=pq\left(1-\frac1p\right)\left(1-\frac1q\right)\\
&=f(pq)
\end{align*}

\bigskip\hrule\bigskip

\textbf{(b)(iii)} $\quad \boxed{\text{False}}$
$$f(1)=1,\qquad f(3)=2$$
$$f(1)f(3)=2=f(3)=f(1\cdot3)$$
But $1$ is not prime

\bigskip\hrule\bigskip

\textbf{(c)} For prime $p$
\begin{align*}
f(p^m)
&=p^m\left(1-\frac1p\right)\\
&=p^{m-1}(p-1)
\end{align*}

$$\because 146410=10\cdot11^4$$
$$\therefore p-1=10,\qquad p^{m-1}=11^4$$
$$\boxed{p=11,\qquad m=5}$$

\bigskip
\hrule

\begin{center}
	This solution was submitted by \textbf{Jon}.\\[4pt]
	Published and refined by \textbf{Jon}.\\[8pt]
	\textit{Last updated: June 20, 2026}
\end{center}

\end{document}